% GENERATED FILE. Edit canonical lessons, then run npm run generate:books.
% canonical-source-hash: fnv1a64:fbc14c06120d0642
% canonical-lessons: MR-C331-vedha, MR-C331-gherne, MR-C331-khandak, MR-C331-tatbandi, MR-C331-buruj

\chapter{The Siege}
\label{ch:mr-the-siege}
\hlchaptermodality{331}

\begin{chapteropening}
\textbf{By the end of this chapter:} \emph{I can describe a fort under siege: the moat, the walls, the towers.}

Complete the last lesson of chapter 331: I can describe a fort under siege: the moat, the walls, the towers.
\end{chapteropening}

\section[\texorpdfstring{\mr{वेढा} (veḍhā)}{वेढा (veḍhā)}]{\mr{वेढा} (veḍhā) — a siege}

\label{lesson:MR-C331-vedha}



\begin{quote}
Before the new one: say the Marathi for to stand guard, then the Marathi for to battle.
\end{quote}

\subsection*{You'll want to know: \mr{वेढा}}
\textbf{\mr{वेढा}} — \emph{veḍhā} — \textquotedblleft{}a siege\textquotedblright{}.

\begin{quote}
\mr{शत्रूचा} \mr{वेढा} \mr{सहा} \mr{महिने} \mr{चालला}.
\end{quote}

\begin{quote}
\emph{śatrūcā veḍhā sahā mahine cāllā.}
\end{quote}

\begin{quote}
The enemy's siege lasted six months.
\end{quote}

Masculine.

\mr{वेढा} \mr{घालणे}, to lay siege; \mr{वेढा} \mr{उठवणे}, to lift it.

\begin{cousinweb}
\textbf{\mr{वेढा}} (\emph{veḍhā}) is inherited, of the family of Sanskrit \emph{veṣṭa}, 'a band, an enclosure': CDIAL lists Marathi \emph{veḍh}, 'a ring, a circumference', under it (CDIAL 12130), and \emph{veḍhṇẽ}, 'to twist, surround', under \emph{veṣṭayati}, 'surrounds' (CDIAL 12132).
\end{cousinweb}

\begin{grammarlens}[title={What you've built}]
One more word for telling what happened.
\end{grammarlens}

\subsection*{Your turn}
\begin{itemize}
\raggedright
  \item \emph{Say it:} \emph{veḍhā}
  \item \emph{Say it:} \emph{veḍhā}, once more
  \item \emph{Say it:} \emph{śatrūcā veḍhā sahā mahine cāllā.} — The enemy's siege lasted six months.
  \item \emph{Recall:} say the Marathi for to stand guard, then the Marathi for to battle, then say \emph{veḍhā} again
\end{itemize}

\begin{tcolorbox}[breakable,colback=teal!4,colframe=teal!35!black,title={Before you move on}]
What does \mr{वेढा} mean? (\textquotedblleft{}A siege\textquotedblright{}.) Say it once more.
\end{tcolorbox}
\section[\texorpdfstring{\mr{घेरणे} (gherṇe)}{घेरणे (gherṇe)}]{\mr{घेरणे} (gherṇe) — to surround}

\label{lesson:MR-C331-gherne}



\begin{quote}
Before the new one: say the Marathi for to battle, then the Marathi for a siege.
\end{quote}

\subsection*{You'll want to know: \mr{घेरणे}}
\textbf{\mr{घेरणे}} — \emph{gherṇe} — \textquotedblleft{}to surround\textquotedblright{}.

\begin{quote}
\mr{सैन्याने} \mr{किल्ल्याला} \mr{चारी} \mr{बाजूंनी} \mr{घेरलं}.
\end{quote}

\begin{quote}
\emph{sainyāne killyālā cārī bājūnnī gherlaṁ.}
\end{quote}

\begin{quote}
The army surrounded the fort on all sides.
\end{quote}

Past, with -\mr{ने} (he …): \mr{त्याने} \mr{घेरलं}, \emph{tyāne gherlaṁ}.

What is surrounded takes -\mr{ला}: \mr{किल्ल्याला} \mr{घेरणे}.

\begin{cousinweb}
\textbf{\mr{घेरणे}} (\emph{gherṇe}) is inherited: CDIAL lists Marathi \emph{gherṇe}, 'to surround', under \emph{*ghir}, 'go round' (CDIAL 4474).
\end{cousinweb}

\begin{grammarlens}[title={What you've built}]
One more word for telling what happened.
\end{grammarlens}

\subsection*{Your turn}
\begin{itemize}
\raggedright
  \item \emph{Say it:} \emph{gherṇe}
  \item \emph{Say it:} \emph{gherṇe}, once more
  \item \emph{Say it:} \emph{sainyāne killyālā cārī bājūnnī gherlaṁ.} — The army surrounded the fort on all sides.
  \item \emph{Recall:} say the Marathi for to battle, then the Marathi for a siege, then say \emph{gherṇe} again
\end{itemize}

\begin{tcolorbox}[breakable,colback=teal!4,colframe=teal!35!black,title={Before you move on}]
What does \mr{घेरणे} mean? (\textquotedblleft{}To surround\textquotedblright{}.) Say it once more.
\end{tcolorbox}
\section[\texorpdfstring{\mr{खंदक} (khandak)}{खंदक (khandak)}]{\mr{खंदक} (khandak) — a moat; a trench}

\label{lesson:MR-C331-khandak}



\begin{quote}
Before the new one: say the Marathi for a siege, then the Marathi for to surround.
\end{quote}

\subsection*{You'll want to know: \mr{खंदक}}
\textbf{\mr{खंदक}} — \emph{khandak} — \textquotedblleft{}a moat; a trench\textquotedblright{}.

\begin{quote}
\mr{किल्ल्याभोवती} \mr{पाण्याने} \mr{भरलेला} \mr{खंदक} \mr{होता}.
\end{quote}

\begin{quote}
\emph{killyābhovtī pāṇyāne bharlelā khandak hotā.}
\end{quote}

\begin{quote}
There was a moat full of water round the fort.
\end{quote}

Masculine.

\mr{खंदक} \mr{खणणे}, to dig a trench.

\begin{cousinweb}
\textbf{\mr{खंदक}} (\emph{khandak}) is a loan from Persian \emph{xandaq}, usually traced to Arabic \emph{khandaq}, 'a trench, a moat'.
\end{cousinweb}

\begin{grammarlens}[title={What you've built}]
One more word for telling what happened.
\end{grammarlens}

\subsection*{Your turn}
\begin{itemize}
\raggedright
  \item \emph{Say it:} \emph{khandak}
  \item \emph{Say it:} \emph{khandak}, once more
  \item \emph{Say it:} \emph{killyābhovtī pāṇyāne bharlelā khandak hotā.} — There was a moat full of water round the fort.
  \item \emph{Recall:} say the Marathi for a siege, then the Marathi for to surround, then say \emph{khandak} again
\end{itemize}

\begin{tcolorbox}[breakable,colback=teal!4,colframe=teal!35!black,title={Before you move on}]
What does \mr{खंदक} mean? (\textquotedblleft{}A moat; a trench\textquotedblright{}.) Say it once more.
\end{tcolorbox}
\section[\texorpdfstring{\mr{तटबंदी} (taṭbandī)}{तटबंदी (taṭbandī)}]{\mr{तटबंदी} (taṭbandī) — a rampart; fortifications}

\label{lesson:MR-C331-tatbandi}



\begin{quote}
Before the new one: say the Marathi for to surround, then the Marathi for a moat.
\end{quote}

\subsection*{You'll want to know: \mr{तटबंदी}}
\textbf{\mr{तटबंदी}} — \emph{taṭbandī} — \textquotedblleft{}a rampart; fortifications\textquotedblright{}.

\begin{quote}
\mr{किल्ल्याची} \mr{तटबंदी} \mr{अजून} \mr{मजबूत} \mr{आहे}.
\end{quote}

\begin{quote}
\emph{killyācī taṭbandī ajūn majbūt āhe.}
\end{quote}

\begin{quote}
The fort's ramparts are still strong.
\end{quote}

Feminine.

\mr{तट}, a wall, + \mr{बंदी}, closing.

\begin{cousinweb}
\textbf{\mr{तटबंदी}} (\emph{taṭbandī}) is built from \emph{taṭ}, 'a wall', from Sanskrit \emph{taṭa}, 'a slope, a bank' (Monier-Williams), and Persian \emph{bandī}, 'a closing, a fastening'.
\end{cousinweb}

\begin{grammarlens}[title={What you've built}]
One more word for telling what happened.
\end{grammarlens}

\subsection*{Your turn}
\begin{itemize}
\raggedright
  \item \emph{Say it:} \emph{taṭbandī}
  \item \emph{Say it:} \emph{taṭbandī}, once more
  \item \emph{Say it:} \emph{killyācī taṭbandī ajūn majbūt āhe.} — The fort's ramparts are still strong.
  \item \emph{Recall:} say the Marathi for to surround, then the Marathi for a moat, then say \emph{taṭbandī} again
\end{itemize}

\begin{tcolorbox}[breakable,colback=teal!4,colframe=teal!35!black,title={Before you move on}]
What does \mr{तटबंदी} mean? (\textquotedblleft{}A rampart; fortifications\textquotedblright{}.) Say it once more.
\end{tcolorbox}
\section[\texorpdfstring{\mr{बुरुज} (buruj)}{बुरुज (buruj)}]{\mr{बुरुज} (buruj) — a bastion; a fort tower}

\label{lesson:MR-C331-buruj}



\begin{quote}
Before the new one: say the Marathi for a moat, then the Marathi for a rampart.
\end{quote}

\subsection*{You'll want to know: \mr{बुरुज}}
\textbf{\mr{बुरुज}} — \emph{buruj} — \textquotedblleft{}a bastion; a fort tower\textquotedblright{}.

\begin{quote}
\mr{बुरुजावरून} \mr{संपूर्ण} \mr{गाव} \mr{दिसतं}.
\end{quote}

\begin{quote}
\emph{burujāvarūn sampūrṇa gāv distaṁ.}
\end{quote}

\begin{quote}
The whole village can be seen from the bastion.
\end{quote}

Masculine.

\mr{बुरुजावर}, on the bastion.

\begin{cousinweb}
\textbf{\mr{बुरुज}} (\emph{buruj}) is a loan from Persian \emph{burj}, usually traced to Arabic \emph{burj}, 'a tower'.
\end{cousinweb}

\begin{grammarlens}[title={What you've built}]
One more word for telling what happened.
\end{grammarlens}

\subsection*{Your turn}
\begin{itemize}
\raggedright
  \item \emph{Say it:} \emph{buruj}
  \item \emph{Say it:} \emph{buruj}, once more
  \item \emph{Say it:} \emph{burujāvarūn sampūrṇa gāv distaṁ.} — The whole village can be seen from the bastion.
  \item \emph{Recall:} say the Marathi for a moat, then the Marathi for a rampart, then say \emph{buruj} again
\end{itemize}

\begin{tcolorbox}[breakable,colback=teal!4,colframe=teal!35!black,title={Before you move on}]
What does \mr{बुरुज} mean? (\textquotedblleft{}A bastion; a fort tower\textquotedblright{}.) Say it once more.
\end{tcolorbox}
